\documentclass[main.tex]{subfiles}
% Import external statements; the generated section retains B28's own numbering.
\begin{luacode*}
  local directory = "registry/e-coordinates/"
  thmlookup.registry_path = directory .. tex.jobname .. ".registry.tsv"
  thmlookup.debug_path = directory .. tex.jobname .. ".debug.log"
  thmlookup.prepare_run()
  for _, predecessor in ipairs(banddeps.graph.B28.predecessors) do
    local base = banddeps.graph[predecessor].artifact_base
    tex.sprint("\\externaldocument{" .. base .. "}[../_" .. predecessor .. ".pdf]")
  end
  tex.sprint("\\externaldocument{" .. directory .. "b28-external}[../_B28.pdf]")
  banddeps.load_predecessor_registries("B28")
  thmlookup.load_registry_file(directory .. "b28-external.registry.tsv")
\end{luacode*}
\renewcommand{\FormulaBandID}{28}
\setcounter{file}{28}
\renewcommand{\theHtable}{coordinates.\arabic{chapter}.\arabic{table}}
\renewcommand{\theHfigure}{coordinates.\arabic{chapter}.\arabic{figure}}

\title{\(e\)-Koordinatenabbildung\\[.5em]\large Definitionen und Beweise zu Band 28}
\author{Martin Kunze}
\date{}
\hypersetup{pdftitle={Band 28 - e-Koordinatenabbildung - Definitionen und Beweise},pdfauthor={Martin Kunze}}
\begin{document}
\hypersetup{pageanchor=false}
\maketitle
\hypersetup{pageanchor=true,hypertexnames=false}
\ProofWideReasonsNearFormula
Diese Ergänzung enthält den vollständigen Abschnitt über Koordinatenisomorphismen
von Halbgruppen aus Band 28 mit dessen ursprünglicher Nummerierung.
Ausgangspunkt sind eine Halbgruppe \((A,\star)\) und ein festes Element \(e\in A\).
Die Abbildung \(x\mapsto(x\star e,e\star x)\), ihre Zielhalbgruppe, die
Multiplikationsabbildung und das Kriterium für einen Koordinatenisomorphismus
werden gemeinsam mit ihren Voraussetzungen und Beweisen dargestellt.
Die Nachschlagetabellen der verwendeten Aussagen erscheinen automatisch
unter den Beweisen.
\par\medskip
Zum Ausgangsband: \FormulaRefAuto{SemigroupDef}[def].
\input{registry/e-coordinates/section.tex}
\end{document}
